\section{Zarathustra}

\begin{quotex}
The eternal Dream is borne on the wings of ageless Light that rends the veil of the vague and goes across Time weaving ceaseless patterns of Being
\flright{\textsc{Rabindranath Tagore}, \emph{Santiniketan}}
\end{quotex}

As initiates from \textbf{Lao Tzu} to \textbf{Fabre d'Olivet} and \textbf{Rene Guenon} show us, there are three powers driving history, the true science of man. \textbf{Destiny} is the resultant of the material, biological, psychic, elemental and astrological forces that act mechanically and by necessity on man. Subject to these forces, man is passive and life is lived through him, in time.

Next, there is \textbf{Will} which is the power of Man to oppose the necessity of Fate; rising above time, he contemplates the flow of events and is the link between Freedom and Necessity. \textbf{Providence} originates from the Divinity and its aim is the perfection of all things. It is free and time does not exist for it, which is just the expression of the movement of eternity.

\paragraph{Background}


From time to time, the action of Providence is revealed in a special and conscious way through the will of a particular man, whether he is called an avatar, bodhisattva, or prophet. He comes at the end of an era to be the spiritual impetus for the next cycle. \textbf{Zarathustra} arrived around 1000 BC or earlier to bring his message to the Persians. The religion of the Persians would have been similar to those of Vedic India and ancient Greece. From the third century BC, knowledge of Zarathustra's teachings became known to the Greeks who considered him a Sage based on his antiquity and the power of his thought. Thus, we can be sure that the ideas of Zarathustra gained currency in the Hellenic and Roman world centuries before the birth of Christ.

\paragraph{Teaching}


In the Hibbert lectures of 1930, published as \textit{The Religion of Man}, \textbf{Rabindranath Tagore} identifies Zarathustra as the Prophet, not of the death of, but rather of the revelation of God; perhaps it is more correct to say that he prophesized the death of the gods. Tagore writes of him:

\begin{quotex}
There can be hardly any question that he was the first man we know who gave a definitely oral character and direction to religion and at the same time preached the doctrine of monotheism which offered an eternal foundation of reality to goodness as an ideal of perfection.
\end{quotex}


The religions of the ancient cities, based on duties, rites, and sacrifices, were degeneration so a new message was necessary. Tagore goes on:

\begin{quotex}
All religions of the primitive type try to keep men bound with regulations of external observances. Zarathustra was the greatest of all the pioneer prophets who showed the path of freedom to man, the freedom of moral choice, the freedom from the blind obedience to unmeaning injunctions, the freedom from the multiplicity of shrines which draw our worship away from the single-minded chastity of devotion.
\end{quotex}


The Prophet sees his teachings as `unheard of words' and a `mystery' because they are a new revelation. Such a revelation is based on knowledge, vidya, or gnosis; its opposite, then, is avidya, ignorance, error. Tagore explains:

\begin{quotex}
The revelation he announce is to him no longer a matter of sentiment, no longer a merely undefined presentiment and conception of the Godhead, but a matter of intellect, of spiritual perception and knowledge ... it is the unbelieving that are unknowing; on the contrary, the believing are learned because they have penetrated into this knowledge.
\end{quotex}


\paragraph{The One God}


\begin{quotex}
When I conceived of Thee, O Mazda, as the very First and the Last, as the most Adorable One, as the Father of the Good Though, as the Creator of Truth and Right, as the Lord Judge of our actions in life, then I made a place for Thee in my very eyes.
\flright{\textsc{Zarathustra}, \emph{Yasna 31.8}}
\end{quotex}


\begin{quotex}
The truth which is not reached through the analytical process of reasoning ... the truth which comes like an inspiration out of context with its surroundings brings with it an assurance that it has been sent from an inner source of divine wisdom, that the individual who has realized it is specially inspired and therefore has his responsibility as a direct medium of communication of Divine Truth.
\flright{\textsc{Tagore}}
\end{quotex}


Thus such ideas became known among the ancients. God could no longer be understood as belonging to the tribe or the city, nor reduced to their `ownmost'. Tagore writes:

\begin{quotex}
The consciousness of God transcends the limitations of race and gathers together all human beings within on spiritual circle of union. Zarathustra was the first prophet who emancipated religion from the exclusive narrowness of the tribal God, the God of a chosen people, and offered it the universal Man. This is a great fact in the history of religion.
\end{quotex}


The Prophet wrote about his enlightenment:

\begin{quotex}
Veriy I believed Thee, O Ahura Mazda, to be the Supreme Benevolent Providence, when Sraosha came to me with the Good Mind, when first I received and became wise with your words. And though the task be difficult, though woe may come to me, I shall proclaim to all mankind Thy message, which Thou declarest to be the best.
\end{quotex}


This \emph{best message} is The Heart of the West\footnote{See Section \ref{sec:HeartWest} in this book.}.


\flrightit{Posted on 2011-07-07 by Cologero}

